\begin{scanpage}{0225}{207}
\noindent is the energy flux density in the incident wave averaged over an oscillation period, and
\[P=\frac{c}{8\pi}\operatorname{Re}\int\mathbf n[\mathbf E\mathbf H^*]dS=\frac12E\sin\vartheta_0\operatorname{Re}\int_{z_1}^{z_2}J(z)e^{ikz\cos\vartheta_0}dz\tag{34.03}\]
is the value of the energy scattered by the vibrator into the surrounding space averaged over a period. Since one may represent the fringing field in the far zone in the direction $\vartheta=\pi-\vartheta_0$ by the equations
\[E_\vartheta=H_\varphi=-\frac{ik}c\sin\vartheta_0\frac{e^{ikR}}R\int_{z_1}^{z_2}J(z)e^{ikz\cos\vartheta_0}dz\tag{34.04}\]
and
\[E_\vartheta=H_\varphi=-E\frac{e^{ikR}}{kR}\cdot F(\pi-\vartheta_0,\vartheta_0),\tag{34.05}\]
then, having determined from this the integral
\[\int_{z_1}^{z_2}J(z)e^{ikz\cos\vartheta_0}dz=\frac{cE}{ik^2\sin\vartheta_0}F(\pi-\vartheta_0,\vartheta_0)\tag{34.06}\]
we obtain
\[S=\frac{\lambda^2}\pi\cos^2\alpha\cdot\operatorname{Im}F(\pi-\vartheta_0,\vartheta_0).\tag{34.07}\]
Calculations of the quantity $S/L^2$ (with $\alpha=0$, $\vartheta_0=\frac\pi2$), performed by us for vibrators with the parameter $\chi=\frac1{2\ln ka}$ taking the values $\chi=-0.05$ and $\chi=-0.1$, are found to be in agreement with the results of Leontovich and Levin [85]. With $\chi=-0.1$, our curve (the dotted line in Figure 76) is only slightly displaced in the direction of longer wavelengths and gives slightly higher resonance peaks.

The effective scattering area $\sigma$ according to the definition equals
\[\sigma(\vartheta,\vartheta_0)=\frac{p_r\cdot4\pi R^2}{p},\tag{34.08}\]
where $p$ is the known quantity (34.02) and
\end{scanpage}
